%======================================================================
\section{Preliminaries}\label{sec:prelim}
%======================================================================

All Hilbert spaces are finite-dimensional. This section fixes the norms in which errors are measured and recalls approximate
representations and their hyperlinear profile, with the two facts that make profiles computable: the defect of a word is at most
linear in its area, and a change of finite presentation changes all defects by bounded factors. It then introduces state-weighted
models, the objects that near-optimal strategies actually produce, and the profiles they define. Finally it recalls linear system
games and fixes the measures of entanglement.

\subsection{Norms and conventions}

Logarithms are to base $2$ unless written $\ln$; $h(x)=-x\log x-(1-x)\log(1-x)$ is the binary entropy, $S$ the von Neumann entropy in
bits, and $D(A\|B)=\Tr A(\log A-\log B)$ the relative entropy of positive operators $A,B$ with $\supp A\subseteq\supp B$. Maps compose
from right to left and groups act on the left. $F(S)$ is the free group on a set $S$, and $\abs{w}$ is the length of a word $w$. For
words or group elements $[a,b]=aba^{-1}b^{-1}$; for operators $[S,T]=ST-TS$ and $\{S,T\}=ST+TS$. An \emph{involution} is a Hermitian
unitary.

For an operator $T$ between Hilbert spaces, $\hs{T}=(\Tr T^*T)^{1/2}$ is the unnormalized Hilbert--Schmidt norm, $\tn{T}$ the trace norm
and $\op{T}$ the operator norm. We use $\hs{XTY}\le\op{X}\hs{T}\op{Y}$, the invariance of $\hs{\cdot}$ under multiplication by unitaries
on either side, and H\"older's inequality $\tn{ST}\le\hs{S}\hs{T}$. On $\C^d$ we write $\tr=\frac1d\Tr$ and $\fn{T}=\hs{T}/\sqrt d$ for the
normalized norm, so $\fn{I}=1$. For unitaries $U,W$ we have $\fn{U-W}=\fn{UW^*-I}$, $\fn{U^{-1}-I}=\fn{U-I}$, and
\begin{equation}\label{eq:unitary-distance}
\fn{U-I}^2=2-2\operatorname{Re}\tr U\le4,\qquad\text{with equality if and only if }U=-I ,
\end{equation}
since $\operatorname{Re}\tr U=-1$ forces every eigenvalue of $U$ to be $-1$. For unitaries $U,V$, $UV-VU=(UVU^{-1}V^{-1}-I)VU$, so
$\fn{UVU^{-1}V^{-1}-I}=\fn{UV-VU}$: inside the normalized norm the group commutator and the operator commutator are interchangeable.
In a fixed state they differ by state commutators (\cref{prelim:models}).

\subsection{Approximate representations}

An approximate representation assigns unitary matrices to the generators of a finite presentation so that every relator is mapped close
to the identity. An element that is nontrivial in the group may still be forced close to the identity, and the hyperlinear profile
measures how large the matrices must be before it can stay away.

\begin{definition}[Approximate representations]\label{def:approx-rep}
Let $S$ be finite, $R\subset F(S)$ finite and $\eps\ge0$. An \emph{$\eps$-representation} of $\langle S\mid R\rangle$ of dimension $d$ is a
homomorphism $\phi:F(S)\to\U(d)$ with $\fn{\phi(r)-I}\le\eps$ for every $r\in R$. We call $\fn{\phi(w)-I}$ the \emph{defect} of $w$.
\end{definition}

\begin{definition}[Hyperlinear profile; {\cite[Definition~2.3]{SV18}}]\label{def:hlp}
For $w\in F(S)$, $0<\delta\le2$ and $\eps>0$, the \emph{hyperlinear profile} $\hlp(w;\delta,\eps)$ is the least $d$ for which some
$\eps$-representation of $\langle S\mid R\rangle$ of dimension $d$ has $\fn{\phi(w)-I}\ge\delta$, and $\infty$ if there is none. We write
$\hlp_R$ when the presentation matters, and say that $w$ is \emph{nontrivial in approximate representations} if
$\hlp(w;\delta,\eps)<\infty$ for some $\delta>0$ and all $\eps>0$.
\end{definition}

The profile is nonincreasing in $\eps$ and nondecreasing in $\delta$. By~\eqref{eq:unitary-distance}, $\delta=2$ asks for $\phi(w)=-I$
exactly; this is the form used by Slofstra and Vidick. If a finitely presented group is hyperlinear, every element $g\ne1$ is nontrivial
in approximate representations, for every finite presentation and every word representing $g$; amenable groups are sofic, hence
hyperlinear~\cite{Gro99,Wei00,ES05,Pes08}.

A word $w$ that represents $1$ in $G=\langle S\mid R\rangle$ is freely equal to a product of $k$ conjugates of elements of $R^{\pm1}$; the
least such $k$ is its \emph{area} $\Area_R(w)$, and the \emph{Dehn function} of the presentation is
$n\mapsto\max\{\Area_R(w):\abs{w}\le n,\ w=1\text{ in }G\}$. For a set of words $W$ we say that $w$ has \emph{area at most $k$ in $W$}
if it is freely equal to a product of $k$ conjugates of elements of $W^{\pm1}$.

\begin{lemma}[Defect is bounded by area]\label{lem:defect}
If $\fn{\phi(r)-I}\le\eps$ for every $r\in W$ and $w$ has area at most $k$ in $W$, then $\fn{\phi(w)-I}\le k\eps$. In particular, in an
$\eps$-representation, $\fn{\phi(w)-\phi(w')}\le\Area_R(w'^{-1}w)\,\eps$ whenever $w,w'$ represent the same element.
\end{lemma}

\begin{proof}
For unitaries $U_1,\dots,U_k$ and $W'=U_2\cdots U_k$, unitary invariance gives $\fn{U_1W'-I}\le\fn{U_1W'-W'}+\fn{W'-I}=\fn{U_1-I}+\fn{W'-I}$,
so $\fn{U_1\cdots U_k-I}\le\sum_i\fn{U_i-I}$ by induction. Write $w$, up to free reduction, as $\prod_{i=1}^ku_ir_i^{\pm1}u_i^{-1}$ with
$r_i\in W$. Since $\phi$ is a homomorphism of the free group, each factor has
$\fn{\phi(u_i)\phi(r_i)^{\pm1}\phi(u_i)^{-1}-I}=\fn{\phi(r_i)^{\pm1}-I}=\fn{\phi(r_i)-I}\le\eps$, and summing gives the first claim. For the
second, $w'^{-1}w$ represents $1$ and $\fn{\phi(w)-\phi(w')}=\fn{\phi(w'^{-1}w)-I}$.
\end{proof}

The lemma is the only way in which approximate representations see the geometry of the group: a short filling forces a word to be nearly
trivial. Its counterpart in a fixed state, where conjugators are no longer free, is \cref{lem:calculus}(d).

\begin{lemma}[Change of presentation]\label{lem:change-presentation}
Let $\langle S\mid R\rangle$ and $\langle S'\mid R'\rangle$ present the same group, fix words $\hat s'\in F(S)$ representing the generators
$s'\in S'$, and let $w\in F(S)$, $w'\in F(S')$ represent the same element. There is $C\ge1$ with
$\hlp_{R'}(w';\delta-C\eps,C\eps)\le\hlp_R(w;\delta,\eps)$ whenever $C\eps<\delta$; if $w'=s'$ is a letter and $\hat s'=w$, the loss in
$\delta$ disappears.
\end{lemma}

\begin{proof}
Let $\vartheta:F(S')\to F(S)$ be the homomorphism $s'\mapsto\hat s'$. The words $\vartheta(r')$, $r'\in R'$, and $\vartheta(w')^{-1}w$
represent $1$; let $C\ge1$ bound their areas in $R$. If $\phi$ is an $\eps$-representation of $\langle S\mid R\rangle$ with
$\fn{\phi(w)-I}\ge\delta$, then $\phi'=\phi\circ\vartheta$ is a $C\eps$-representation of $\langle S'\mid R'\rangle$ of the same dimension
by \cref{lem:defect}, and $\fn{\phi'(w')-\phi(w)}\le C\eps$, so $\fn{\phi'(w')-I}\ge\delta-C\eps$. In the special case
$\phi'(w')=\phi(\hat s')=\phi(w)$.
\end{proof}

So lower bounds of the form $\exp(c\,(\delta/\eps)^{\alpha})$, and finiteness of the profile, pass between finite presentations, with new
constants. The same holds for the state-weighted profiles below (\cref{state:change-presentation}).

\subsection{State-weighted models}\label{prelim:models}

Let $\rho$ be a density matrix on $\C^d$ and $\xi=\rho^{1/2}$, so that $\hs{\xi}=1$. For an operator $Y$ on $\C^d$ we call
\[
\hs{Y\xi}=\Tr(Y^*Y\rho)^{1/2}\qquad\text{and}\qquad\hs{[Y,\xi]}
\]
the \emph{state-weighted norm} of $Y$ and the \emph{state commutator} of $Y$. For the \emph{flat state} $\rho=I/d$ the first is $\fn{Y}$
and the second vanishes. We use without comment that $\hs{XY\xi}\le\op{X}\hs{Y\xi}$, that $\hs{Y\xi U}=\hs{Y\xi}$ for unitary $U$, that
$\hs{\xi Y}=\hs{Y^*\xi}$, that $\hs{(U^*-I)\xi}=\hs{(U-I)\xi}$ for unitary $U$, and that $[U^*,\xi]=-U^*[U,\xi]U^*$, so
$\hs{[U^*,\xi]}=\hs{[U,\xi]}$. For contractions $S,T$ the Leibniz rule $[ST,\xi]=S[T,\xi]+[S,\xi]T$ gives
$\hs{[ST,\xi]}\le\hs{[S,\xi]}+\hs{[T,\xi]}$. For unitaries $U,V$ the identities $[U,V]=(UVU^{-1}V^{-1}-I)VU$ and
$UVU^{-1}V^{-1}-I=[U,V]U^{-1}V^{-1}$ show that $\hs{[U,V]\xi}$ and $\hs{(UVU^{-1}V^{-1}-I)\xi}$ differ by at most
$2\hs{[U,\xi]}+2\hs{[V,\xi]}$.

\begin{definition}[State-weighted models]\label{def:model}
Let $\langle S\mid R\rangle$ be a finite presentation, $w\in F(S)$ and $e\ge0$. A \emph{state-weighted model} of $\langle S\mid R\rangle$
at precision $e$ on $\C^d$ is a pair $(\phi,\rho)$ of a homomorphism $\phi:F(S)\to\U(d)$ and a density matrix $\rho$ on $\C^d$ such that,
with $\xi=\rho^{1/2}$,
\begin{enumerate}[label=(S\arabic*),leftmargin=3.2em]
\item\label{prelim:S1} $\hs{(\phi(r)-I)\xi}\le e$ for every $r\in R$;
\item\label{prelim:S2} $\hs{[\phi(s),\xi]}\le e$ for every $s\in S$.
\end{enumerate}
It is a model of the pair $(\langle S\mid R\rangle,w)$ if moreover
\begin{enumerate}[label=(S\arabic*),leftmargin=3.2em,resume]
\item\label{prelim:S3} $\hs{(\phi(w)+I)\xi}\le e$.
\end{enumerate}
We call $\hs{(\phi(u)-I)\xi}$ the \emph{defect of $u$ in the state}.
\end{definition}

Condition~\ref{prelim:S2} has no flat analogue: a trace ignores conjugation, a state does not, and~\ref{prelim:S2} is what pays for conjugators
(\cref{lem:calculus}). Condition~\ref{prelim:S3} is the state-weighted form of $\phi(w)=-I$.

\begin{definition}[State-weighted profiles]\label{def:state-profile}
Let $\langle S\mid R\rangle$ be a finite presentation, $w\in F(S)$ and $e>0$. The \emph{state-weighted profile} $\hlps(w;e)$ is the least $d$
for which a model of $(\langle S\mid R\rangle,w)$ at precision $e$ on $\C^d$ exists, that is, a pair $(\phi,\rho)$
satisfying~\ref{prelim:S1}--\ref{prelim:S3}; it is $\infty$ if there is none. The \emph{entropy profile} $\hS(e)$ is the infimum of $S(\rho)$ over all such
models, in any dimension. We write $\hlps_R$ and $\hS_R$ when the presentation matters.
\end{definition}

Both profiles are nonincreasing in $e$, and $\hS(e)\le\log\hlps(w;e)$ since $S(\rho)\le\log d$. If $\phi$ is an $e$-representation with
$\fn{\phi(w)-I}\ge2$, that is with $\phi(w)=-I$, then $(\phi,I/d)$ satisfies~\ref{prelim:S1}--\ref{prelim:S3} at precision $e$. Hence
\[
\hlps(w;e)\le\hlp(w;2,e),
\]
and a lower bound on $\hlps$ is a stronger statement than the same lower bound on $\hlp(\,\cdot\,;2,\cdot)$.

\subsection{Linear system games}

Let $A\in\F_2^{m\times n}$ and $b\in\F_2^m$, and let $V_i=\{j:A_{ij}=1\}$ be the set of variables of the $i$-th equation, assumed nonempty;
put $s=\max_i\abs{V_i}$. We write answers as signs, the bit $0$ as $+1$ and the bit $1$ as $-1$.

\begin{definition}[Linear system games; \cite{CM14,CLS17,SV18}]\label{def:lcs}
\leavevmode
\begin{enumerate}[label=(\alph*),itemsep=1pt]
\item The \emph{solution group} $\Gamma=\Gamma(A,b)$ is generated by $x_1,\dots,x_n,J$ with the relators
\begin{gather*}
x_j^2,\quad J^2,\quad [x_j,J]\ \ (1\le j\le n),\qquad [x_j,x_k]\ \ (j,k\in V_i\text{ for some }i),\\
J^{b_i}x_{j_1}x_{j_2}\cdots x_{j_r}\ \ (1\le i\le m,\ V_i=\{j_1<\dots<j_r\}).
\end{gather*}
The element $J$ is a central involution. Approximate representations and models of $\Gamma$ are taken with respect to this presentation,
and with $w=J$; we write $\hlp_\Gamma$, $\hlps_\Gamma$ and $\hS_\Gamma$.
\item In the game $G=G(A,b)$ with \emph{question distribution} $\pi$ on $\{1,\dots,m\}\times\{1,\dots,n\}$, a referee draws $(i,j)$ from
$\pi$. Alice receives $i$ and answers $a\in\{\pm1\}^{V_i}$ with $\prod_{j\in V_i}a_j=(-1)^{b_i}$; Bob receives $j$ and answers
$c\in\{\pm1\}$. They win if $j\notin V_i$ or $a_j=c$. We put $\pi_{\min}=\min\{\pi(i,j):j\in V_i\}$ and assume $\pi_{\min}>0$. The
\emph{default distribution} draws $i$ uniformly and then $j\in V_i$ uniformly, so $\pi_{\min}\ge1/(ms)$; Slofstra and Vidick draw $i$
and $j$ independently and uniformly~\cite{SV18}, so $\pi_{\min}=1/(mn)$.
\item A \emph{strategy} consists of a unit vector $\psi\in H_A\otimes H_B$, a POVM $\{P^i_a\}_a$ on $H_A$ for each $i$, indexed by the
answers to $i$, and a POVM $\{Q^j_c\}_{c=\pm1}$ on $H_B$ for each $j$. Its \emph{observables} are the Hermitian contractions
$A_{ij}=\sum_aa_jP^i_a$ ($j\in V_i$) and $B_j=Q^j_{+1}-Q^j_{-1}$, and the pair $(i,j)$ with $j\in V_i$ is won with probability
\[
\sum_a\bigl\langle\psi,(P^i_a\otimes Q^j_{a_j})\psi\bigr\rangle=\tfrac12\bigl(1+\langle\psi,(A_{ij}\otimes B_j)\psi\rangle\bigr).
\]
The winning probability $\omega(G;\mathcal S)$ of a strategy $\mathcal S$ is the $\pi$-average of these probabilities, a pair with
$j\notin V_i$ counting as won. The \emph{quantum value} $\omq(G)$ is the supremum of $\omega(G;\mathcal S)$ over all strategies, and
$\omqc(G)$ the supremum over commuting-operator strategies, in which $\psi$ lies in one Hilbert space and Alice's operators commute with
Bob's.
\end{enumerate}
\end{definition}

The following facts are standard. A perfect commuting-operator strategy exists if and only if $J\ne1$ in $\Gamma$~\cite[Theorem~4]{CLS17}.
A perfect strategy on a tensor product of Hilbert spaces, of finite or infinite dimension, exists if and only if some finite-dimensional
representation of $\Gamma$ sends $J$ to $-I$~\cite[Theorem~1]{CM14}, \cite[Theorem~5]{CLS17}, \cite[Theorem~3.2]{Slo19}; a
finite-dimensional representation in which $J\ne1$ restricts to the $(-1)$-eigenspace of the image of $J$, where $J$ acts as $-I$. Finally, $\omq(G)=1$ if
and only if $J$ is nontrivial in approximate representations of $\Gamma$: the ``if'' direction is~\cite[Proposition~3.4]{Slo19}
(restated as~\cite[Proposition~4.3]{SV18}), and the ``only if'' direction follows from the lower bound
$\hlp_\Gamma(J;2,C\eps^{1/4})\le\Ent(G,\eps)$ of~\cite[Corollary~5.2]{SV18}, and~\cite[Corollary~5.3]{SV18} states the equivalence. Whether $\omq(G)=1$ does not depend on $\pi$, because a
strategy that wins with probability $1-\eps$ loses each pair $(i,j)$ with $j\in V_i$ with probability at most $\eps/\pi_{\min}$.
\Cref{sec:transfer} proves quantitative and state-weighted forms of these correspondences.

\subsection{Entanglement}

A unit vector $\psi\in H_A\otimes H_B$ has a Schmidt decomposition $\psi=\sum_{k=1}^r\sqrt{\lambda_k}\,e_k\otimes f_k$ with $\lambda_k>0$
and orthonormal families $(e_k)$, $(f_k)$. Its reduced density matrices $\rho_A$ and $\rho_B$ have the same nonzero eigenvalues
$\lambda_k$. We keep apart the \emph{Schmidt rank} $r=\rank\rho_B$, which is at most both local dimensions, and the \emph{entanglement
entropy} $S(\rho_A)=S(\rho_B)\le\log r$, the number of ebits of $\psi$. For a mixed state $\sigma$ on $H_A\otimes H_B$ the
\emph{entanglement of formation} is~\cite{BDSW96}
\[
E_F(\sigma)=\inf\Bigl\{\sum_jp_j\,S\bigl(\Tr_B|\psi_j\rangle\langle\psi_j|\bigr):\ \sigma=\sum_jp_j|\psi_j\rangle\langle\psi_j|\Bigr\}.
\]
The distillable entanglement is at most $E_F$, so lower bounds on $E_F$ say nothing about it.

\begin{definition}[Dimension needed to play; \cite{SV18}]\label{def:entanglement}
For a nonlocal game $G$ and $\eps>0$, $\Ent(G,\eps)$ is the least $d$ such that some strategy with $H_A=H_B=\C^d$ wins with probability at
least $\omq(G)-\eps$, and $\Entme(G,\eps)$ is the least $d$ for which such a strategy exists with the maximally entangled state
$\Phi_d=d^{-1/2}\sum_ke_k\otimes e_k$ of $\C^d\otimes\C^d$. So $\Ent\le\Entme$, and when $\omq(G)=1$ both count strategies that win with
probability at least $1-\eps$.
\end{definition}

\begin{remark}[Local dimension, Schmidt rank and projective measurements]\label{rem:local-dimension}
Two normalizations are free. \emph{Schmidt-rank compression:} if $\psi=\sum_{k\le r}\sqrt{\lambda_k}e_k\otimes f_k$ and $\Pi_A,\Pi_B$
project onto the spans of the $e_k$ and of the $f_k$, then compressing every $P^i_a$ by $\Pi_A$ and every $Q^j_c$ by $\Pi_B$ gives POVMs
on two spaces of dimension $r$ with the same probabilities, since $\psi=(\Pi_A\otimes\Pi_B)\psi$. Hence $\Ent(G,\eps)$ is also the least
Schmidt rank, and the least $\min(\dim H_A,\dim H_B)$, of a strategy that wins with probability at least $\omq(G)-\eps$; a lower bound on
the Schmidt rank of every such strategy is a lower bound on $\Ent(G,\eps)$, and a lower bound on its entanglement entropy is a lower bound
on $\log\Ent(G,\eps)$. \emph{Naimark dilation:} if $K$ is the largest number of answers to a question, dilating each of Alice's POVMs to a
projective measurement on $H_A\otimes\C^K$, with the state $\psi\otimes|0\rangle$, changes neither Bob's reduced density matrix nor any
probability, and multiplies Alice's local dimension by at most $K$.
\end{remark}
